\section{Three ways to be right for the wrong reason}
\label{sec:failures}

\paragraph{(i) The audio gives the answer away.}
\label{sec:conf}
The decoder sees the candidate audio, so anything in the audio that marks the
attended talker solves the task without the brain. Holding out listeners or
stimuli does not help, because the mark belongs to the \emph{role} of attended
talker, so it is present for every listener and every stimulus. Making all
candidates equally loud does not help either. Loudness is removed exactly by
standardising each candidate, but standardising only fixes the mean and the
variance. Every statistic that does not change with scale survives: skew,
kurtosis, sparsity and the balance between frequency bands. On this dataset the
attended and ignored talkers differ in exactly these. The peak-to-average ratio of
the envelope alone separates them with an effect size of
$d=\R{sep/envelopecrestdB/d}$. The dataset's authors trace this to the loudness
adjustment, which clipped the attended talker's peaks \citep{maestro2026}.

We therefore test each candidate set \emph{before} training, with an
\textbf{acceptance probe}: a logistic regression on eight statistics that do not
depend on scale, evaluated on held-out content. A clean candidate set must score
$1/K$. The \cUnm{} set scores \R{probe/raw/K4/excess} above chance.
\cMat{} brings this down to \R{probe/qmatch/K4} against $0.25$, and
same-talker candidates to \R{probe/shifted_qm/K2} against $0.50$
(Table~\ref{tab:probe}).

\paragraph{(ii) The standard scoring function has a shortcut.}
\label{sec:degenerate}
The standard score normalises the brain and audio embeddings at each time step,
multiplies them, averages over time and applies a linear layer
\citep{accou2021,monesi2020}. Suppose the brain encoder outputs the same unit
vector $\hat c$ at every step. This is not a hypothetical state: once the audio
already carries the answer, the easiest way for gradient descent to lower the
loss is to let the encoder ignore its noisy input and emit a constant, and we
observe exactly this collapse when the standard model is trained on this data
(Section~\ref{sec:res-ladder}). With a constant $\hat c$ the time average splits
apart and the score becomes $w^\top(\hat c \odot \frac1T\sum_t \hat a_{k,t}) + b$,
where $\hat a_{k,t}$ is the audio embedding of candidate $k$ at time step $t$,
$T$ is the number of time steps in the window, $\odot$ is the elementwise
product, and $w$ and $b$ are the weights and bias of the final linear layer.
This is a linear classifier on the time-averaged audio, with exactly the power
needed to read the audio mark of (i). So the architecture contains an audio-only
decoder, which it reaches by letting the brain encoder go constant. Problem (i)
makes the audio shortcut \emph{enough}; problem (ii) makes it \emph{reachable}.

\paragraph{(iii) Easy streams crowd out the brain.}
\label{sec:freeload}
When several streams are trained together, the model relies most on the easiest
one \citep{wang2020gradblend,huang2022competition,wu2022greedy,peng2022ogm}.
Here this changes what the result means. Gaze, head motion and scene video show
where the listener is \emph{facing}, which in a room of loudspeakers predicts the
attended talker well \citep{kidd2005,rotaru2024}. Only the EEG follows what the
talker is \emph{saying}. A fused decoder that quietly relies on gaze reports a
higher number while supporting a weaker claim.

As Section~\ref{sec:intro} argues, it also fails exactly for the listeners a brain-based device
is for: those who cannot, or choose not to, face the talker they listen to.

\paragraph{One test catches all three.}
\label{sec:perm}
Shuffling a recorded stream across test windows, while each window keeps its own
candidates and label, costs accuracy only if the decoder uses that stream. We
shuffle the recording, not the labels: shuffling the labels would also break the
link between the audio and the answer, and would hide the very shortcut we are
testing for (Appendix~\ref{app:perm}). Let
$v(T)$ be the accuracy when the streams in the set $T$ are left intact and the
others are shuffled. Then $v(\emptyset)$ is what \emph{this} model can get from the
audio alone, measured through its own audio pathway. The probe bounds what a
simple reader can take from the audio; the shuffled accuracy measures what this
decoder actually takes. Appendix~\ref{app:failures} gives all three problems in
more detail.
